\scanpage{185}
avec donc
\[
(1\wedge x\wedge y\wedge z)\otimes e^{-1}=2\varepsilon_0
\]
\margin{$\varepsilon_0=\bigl[\alpha\wedge\alpha'\wedge x_0\wedge y_0\bigr)\unsure{e^{-1}}\in\pm1$}
le signe dépend de la convention d'orient. (qu'il faudra préciser le moment venu !).

D'autre part
\[
\deleted{\operatorname{Tr}xyz=1}\qquad
\tau-\tau'=\operatorname{Tr}\underbrace{[x,y]}_z.z=\operatorname{Tr}z^2=\operatorname{Tr}1=2
\]
d'où
\[
2=2\varepsilon\varepsilon_0\qquad\text{i.e.}
\]
\[
\varepsilon=\deleted{\unsure{1}}\inserted{\unsure{\tfrac12}}\varepsilon_0\quad\deleted{\unsure{\varepsilon}}
\]
d'où
\[
\boxed{\tau-\tau'=\deleted{\unsure{1}}\inserted{\unsure{\tfrac12}}\varepsilon_0\underbrace{(1\wedge x\wedge y\wedge z)\otimes e^{-1}}_{\Delta\ \mathrm{def}}}
\quad\text{i.e. }2(\tau-\tau')=\varepsilon_0\Delta
\]
(a un sens si $2$ inv. dans l'anneau de base $k$ où on travaille).

On en tire
\[
4\tau\tau'=(\tau+\tau')^2-(\tau-\tau')^2=\Sigma^2-\deleted{\unsure{1}}\inserted{\unsure{\tfrac14}}\Delta^2
\]
À vérifier (dans $\mathbb{Q}[[\deleted{b,b',b''}\ x_{ij},y_{ij},z_{ij}]]$ $1\leq i,j\leq2$)
\[
\Sigma^2-\deleted{\unsure{1}}\inserted{\unsure{\tfrac14}}\Delta^2\in4\mathbb{Z}[b,b',b'',t,t',t'']
\]
\margin{?}
\deleted{i.e.}
\deleted{\[
4\Sigma^2-\Delta^2\in16\mathbb{Z}[b,b',b'',t,t',t'']
\]}
d'où l'expression
\[
\tau\tau'=\Pi(b,b',b'',t,t',t'')\in\mathbb{Z}[b,b',b'',t,t',t'']
\]
expression polynomiale de poids total $6$ (les $b$ étant de
\margin{184}
